\section{Coincidence}
\label{sec:coincidence}

Three things you can do this afternoon.

\paragraph{Two rulers.}
Take a ruler with ten divisions to the inch and lay it edge to edge against
one with nine. Somewhere along the pair a tick of the first sits exactly on a
tick of the second. Move your eye along. The next pair of ticks is a hair
apart, the next a hair more, and after ten ticks of the first ruler and nine
of the second the ticks coincide again. What you see, without trying, is a
slow pattern of near-coincidences that repeats every inch. It is on neither
ruler. Nothing about it needs to be exact, either: bend the rulers a little,
or space the ticks slightly unevenly, and the pattern of coincidences is
still there, slightly uneven in its own way.

\paragraph{Two click trains.}
Now sound. Two trains of clicks at nearly equal rates, say a hundred and a
hundred and one per second, go into a detector that fires only when two
clicks arrive together. Clicks arrive together once per second, so the
detector fires once per second. That is the difference of the two rates,
and it is all the detector knows. Raise both trains to a thousand and a
thousand and one clicks per second and the detector's output is
\emph{exactly} the same, one click per second; it never learned the
original rates at all. The sum of the two trains, as a signal, contains
nothing that repeats once per second. The detector makes it.

\paragraph{Two gratings.}
Print a grating of fine lines on a transparency, print another with a
slightly different spacing, and lay one over the other. From arm's length
you see broad bands, much wider than the lines. Blur your eyes: the bands
stay, and the lines go. Photograph the pair out of focus, threshold the
photograph to pure black and white, print it badly: the bands are always in
the same place. Only their contrast changes.

\paragraph{What the three share.}
Each family comes in numbered members: ticks, clicks, lines. The pattern
that emerges is a pattern of \emph{coincidences} between the numbering of
one family and the numbering of the other. It lives on a scale set by how
slowly the two numberings drift apart, and it is read out by anything that
\emph{pools} --- an eye, a detector, a blur --- with the pooling changing
only how strongly the pattern shows, never where it is. None of the three
descriptions used the word frequency, and the second example does not even
have a stable one. What the three have is a count.

\begin{definition}[The count]
\label{def:count}
A \emph{family} is anything that comes in numbered members along a line, a
plane, a circle or a clock: ticks, clicks, lines, rings, teeth, steps, days.
Its \emph{count} $\cnt(x)$ is the number of members up to the point $x$,
interpolated between members, so that $\cnt$ is a whole number exactly on a
member and $\cnt=n+\tfrac12$ halfway from member $n$ to member $n+1$.
\end{definition}

For a ruler of pitch $T$ the count is $\cnt(x)=x/T$: ten members to the inch
means $\cnt$ rises by ten per inch. But the count exists for any family,
evenly spaced or not, and that is the whole reason this paper never needs a
frequency. A frequency says where a family is on average. A count says where
it is here. Figure~\ref{fig:rulers} draws both rulers with their counts and
the one number that the pattern of coincidences depends on: the difference
of the two counts.

\begin{figure}[t]
  \centering
  \begin{tikzpicture}[x=1cm, y=1cm, font=\scriptsize, >={Stealth[length=3pt]}]
    % Bands where the ticks nearly coincide: the third pattern.
    \foreach \c in {0, 5.5, 11} {
      \fill[accent!12] ({\c-0.9},-0.05) rectangle ({\c+0.9},2.45);
    }
    % Ruler A: pitch 0.5 (10 per "inch" of 5 units): ticks at 0.5 k.
    \draw[ink, line width=0.6pt] (0,1.6) -- (11,1.6);
    \foreach \k in {0,...,22} {
      \draw[ink, line width=0.5pt] ({0.5*\k},1.6) -- ({0.5*\k},2.1);
      \node[ink, above=0pt] at ({0.5*\k},2.1) {\k};
    }
    \node[ink, anchor=east] at (-0.15,1.85) {family A};
    % Ruler B: pitch 0.55 (9 per 5 units): ticks at 0.55 k.
    \draw[ink, line width=0.6pt] (0,0.9) -- (11,0.9);
    \foreach \k in {0,...,20} {
      \draw[ink, line width=0.5pt] ({0.55*\k},0.9) -- ({0.55*\k},0.4);
      \node[ink, below=0pt] at ({0.55*\k},0.4) {\k};
    }
    \node[ink, anchor=east] at (-0.15,0.65) {family B};
    % One period of the pattern.
    \draw[accent, ->, line width=0.6pt] (0,2.75) -- (5.5,2.75) node[midway, above, accent] {one period of the third pattern};
    \draw[accent, ->, line width=0.6pt] (5.5,2.75) -- (0,2.75);
    % The counts, and their difference, as small plots below.
    \begin{scope}[shift={(0,-2.9)}]
      \draw[soft, ->] (0,0) -- (11.3,0) node[right, ink] {$x$};
      \draw[soft, ->] (0,0) -- (0,2.4) node[above, ink] {counts};
      \draw[cool, line width=0.8pt] (0,0) -- (11,2.2) node[right, cool] {$\cnt_A=x/0.5$};
      \draw[cool!60, line width=0.8pt, dashed] (0,0) -- (11,2.0) node[right, cool!70, yshift=-6pt] {$\cnt_B=x/0.55$};
      \draw[accent, line width=1pt] (0,0) -- (11,0.4) node[right, accent] {$D=\cnt_A-\cnt_B$};
      \foreach \c/\v in {5.5/1, 11/2} {
        \draw[accent, dotted] (\c,0) -- (\c,{0.2*\v});
        \node[accent, below] at (\c,0) {$D=\v$};
      }
      \node[accent, below] at (0,0) {$D=0$};
    \end{scope}
  \end{tikzpicture}
  \caption{\figlead{Two rulers and the one number that matters.} Family A
  has ten members where family B has nine. Above: the rulers, with their
  members numbered; the shaded bands are where a member of A falls near a
  member of B, the pattern that is on neither ruler. Below: the two counts
  $\cnt_A$ and $\cnt_B$ rise at slightly different rates, and their
  difference $D=\cnt_A-\cnt_B$ rises slowly. The bands sit exactly where $D$
  is a whole number. The pattern repeats when $D$ has grown by one, which
  takes ten members of A and nine of B: the drawing is the definition of a
  vernier, and everything in this paper is a generalisation of it.}
  \label{fig:rulers}
\end{figure}

\paragraph{What this paper claims.}
Every superposition of families is a point moving on a torus, plus a fixed
picture painted on the torus (\S\ref{sec:torus}). Any observer that pools
before it decides --- a blur, an exposure, an ear, a detector --- sees the
picture averaged along the torus's fast directions, and that average is a
function of the \emph{slow combinations of counts} alone; it has a closed
form, and the observer's own nonlinearity changes the picture but never the
geometry (\S\ref{sec:pooling}). Which of the possible combinations shows is
a question of arithmetic with a classical answer and a twist
(\S\ref{sec:selection}). The pattern that emerges has its own count, so it
can beat again, and the hierarchy is exact (\S\ref{sec:iterated}). Counts
can fail to exist in exactly three ways, and each way has a signature
(\S\ref{sec:counting}). Everything with a number attached is measured
(\S\ref{sec:ledger}), most of it in a drawing tool built for the purpose
(\S\ref{sec:instrument}), and the predictions that reach outside drawings
--- into sound, into strobes, into anything that counts --- are the point.

\begin{tryit}{the rulers}
Two rulers with different units work (centimetres against sixteenths of an
inch is a good pair), or two combs, or two pieces of window screen laid
over each other and turned slowly. Count how many members of each family fit
in one period of the bands. In Figure~\ref{fig:rulers} that is ten and nine;
the rule you will find is that the two counts differ by exactly one per
period, whatever the families are.
\end{tryit}
